\chapter*{Lời Mở Đầu}
\addcontentsline{toc}{chapter}{Lời Mở Đầu}
\markboth{Lời Mở Đầu}{Lời Mở Đầu}

\begin{truyen}{Lời Mở Đầu}{Opening Words}
\chuhoa{T}{riết lý nhân sinh không chỉ là của các nhà triết học — nó ở ngay trong những chuyện đời thường.}
\textit{(Life philosophy is not just for philosophers --- it lives right inside everyday stories.)}

Quyển bốn đưa người đọc đến với những tình huống quen thuộc từ cổ chí kim: một người thầy thay đổi học trò, một lời khuyên đúng lúc thay đổi số phận, một bài học từ thất bại mang lại thành công.
\textit{(Volume four takes the reader through familiar situations from ancient to modern times: a teacher changing a student, timely advice changing a fate, a lesson from failure leading to success.)}

Triết lý nhân sinh không phải để tranh luận — nó để sống. Mỗi câu chuyện trong quyển này là một minh chứng cho điều đó.
\textit{(Life philosophy is not for debate --- it is for living. Every story in this volume is a testament to that.)}

Đọc, suy ngẫm và áp dụng — đó là cách tốt nhất để học triết lý nhân sinh.
\textit{(Read, reflect, and apply --- that is the best way to learn life philosophy.)}
\end{truyen}

\begin{baihoc}
Người khôn không chỉ học từ kinh nghiệm của mình --- họ còn học từ lịch sử và câu chuyện của người khác.
\textit{(Wise people do not only learn from their own experience --- they also learn from history and others' stories.)}
\end{baihoc}

\begin{tcolorbox}[colback=truyenblue!6,colframe=truyenblue,coltitle=white,fonttitle=\bfseries,title={Easy English},arc=4pt,boxrule=1pt,breakable]
\textbf{Short English to remember:}

Philosophy lives in stories, not lectures.

Ancient wisdom speaks to modern problems.
\end{tcolorbox}

\begin{tcolorbox}[colback=truyengold!10,colframe=truyengold!70!black,coltitle=black,fonttitle=\bfseries,title={Tự Học Nhanh},arc=4pt,boxrule=1pt,breakable]
1. Triết lý nào trong quyển này em thấy mình cần nghĩ lại? \textit{Which philosophy in this volume do you feel you need to reconsider?}

2. Lịch sử nào trong sách nhắc em nhớ đến một ai đó trong cuộc sống? \textit{Which historical figure in the book reminds you of someone in your life?}

3. Em có thể rút ra bài học nào từ thất bại của một nhân vật trong sách? \textit{What lesson can you draw from a failure of a character in the book?}
\end{tcolorbox}

\ngancach
